% ============================================================================
% preamble.tex — 包、字体、排版、代码高亮、自定义命令
% ============================================================================

% --- 中文支持 (XeLaTeX) ---
\usepackage{ctex}              % 中文排版核心
\usepackage{fontspec}          % 字体选择

% --- 页面布局 ---
\usepackage[
  a4paper,
  top=2.5cm, bottom=2.5cm,
  left=2.5cm, right=2.5cm
]{geometry}

% --- 数学 ---
\usepackage{amsmath, amssymb}

% --- 图表 ---
\usepackage{graphicx}
\usepackage{float}
\usepackage{booktabs}          % 专业表格线
\usepackage{longtable}
\usepackage{tabularx}

% --- 绘图（内存布局、数据结构、bitfield 图、指针关系） ---
\usepackage{tikz}
\usetikzlibrary{
  positioning,    % 相对定位 (right=of, below=of)
  calc,           % 坐标计算
  arrows.meta,    % 箭头样式
  decorations.pathreplacing,  % 花括号标注
  fit,            % 包围框
  shapes.multipart,  % 多段节点（struct field 图）
}

% 内存块样式
\tikzset{
  memcell/.style={
    draw, minimum width=2cm, minimum height=0.6cm,
    font=\ttfamily\small, anchor=west,
  },
  memaddr/.style={
    font=\ttfamily\scriptsize\color{gray}, anchor=east,
  },
  pointer/.style={
    -{Stealth[length=4pt]}, thick, blue!60!black,
  },
  bitcell/.style={
    draw, minimum width=0.7cm, minimum height=0.6cm,
    font=\ttfamily\small, anchor=west,
  },
  fieldlabel/.style={
    font=\small\itshape, text=blue!50!black,
  },
}

% --- 颜色与代码高亮 ---
\usepackage{xcolor}
\usepackage{listings}

% 代码样式：C 语言
\lstdefinestyle{cstyle}{
  language=C,
  basicstyle=\ttfamily\small,
  keywordstyle=\color{blue}\bfseries,
  commentstyle=\color{gray}\itshape,
  stringstyle=\color{red!70!black},
  numbers=left,
  numberstyle=\tiny\color{gray},
  numbersep=8pt,
  frame=single,
  framesep=4pt,
  breaklines=true,
  tabsize=4,
  showstringspaces=false,
  columns=flexible,
}

% 代码样式：x86 汇编（NASM）
\lstdefinestyle{nasmstyle}{
  language=[x86masm]Assembler,
  basicstyle=\ttfamily\small,
  keywordstyle=\color{blue}\bfseries,
  commentstyle=\color{gray}\itshape,
  numbers=left,
  numberstyle=\tiny\color{gray},
  numbersep=8pt,
  frame=single,
  framesep=4pt,
  breaklines=true,
  tabsize=4,
  showstringspaces=false,
  columns=flexible,
}

% 代码样式：Shell / Makefile
\lstdefinestyle{shellstyle}{
  language=bash,
  basicstyle=\ttfamily\small,
  keywordstyle=\color{blue},
  commentstyle=\color{gray}\itshape,
  frame=single,
  framesep=4pt,
  breaklines=true,
  tabsize=4,
  showstringspaces=false,
  columns=flexible,
}

\lstset{style=cstyle}  % 默认 C

% --- 超链接 ---
\usepackage{hyperref}
\hypersetup{
  colorlinks=true,
  linkcolor=blue!60!black,
  urlcolor=blue!80!black,
  citecolor=green!50!black,
  pdfauthor={SZY},
  pdftitle={自己动手写操作系统},
}

% --- 页眉页脚 ---
\usepackage{fancyhdr}
\pagestyle{fancy}
\fancyhf{}
\fancyhead[L]{\leftmark}
\fancyhead[R]{\thepage}
\renewcommand{\headrulewidth}{0.4pt}

% --- 自定义命令 ---

% 寄存器名
\newcommand{\reg}[1]{\texttt{#1}}

% 十六进制地址
\newcommand{\hex}[1]{\texttt{0x#1}}

% 关键术语首次出现
\newcommand{\term}[1]{\textbf{#1}}

% Bitfield 描述环境
\newenvironment{bitfield}
  {\begin{description}\setlength{\itemsep}{0pt}}
  {\end{description}}
\newcommand{\bit}[2]{\item[\texttt{bit #1}] #2}

% 设计决策高亮框
\usepackage{tcolorbox}
\tcbuselibrary{breakable, skins}

\newtcolorbox{designbox}[1][]{
  colback=blue!3,
  colframe=blue!40!black,
  fonttitle=\bfseries,
  title={设计决策},
  breakable,
  #1
}

\newtcolorbox{whybox}[1][]{
  colback=green!3,
  colframe=green!40!black,
  fonttitle=\bfseries,
  title={为什么？},
  breakable,
  #1
}

\newtcolorbox{warnbox}[1][]{
  colback=red!3,
  colframe=red!40!black,
  fonttitle=\bfseries,
  title={注意},
  breakable,
  #1
}

% CPU 状态说明框
\newtcolorbox{cpustate}[1][]{
  colback=orange!3,
  colframe=orange!40!black,
  fonttitle=\bfseries,
  title={CPU 状态变化},
  breakable,
  #1
}

% --- 思考题与练习 ---

% 计数器（按章重置）
\newcounter{thinknum}[chapter]
\newcounter{exercisenum}[chapter]

% 思考题（穿插在正文中的小问题，引导读者停下来想一想）
\newtcolorbox{thinkbox}[1][]{
  colback=purple!3,
  colframe=purple!40!black,
  fonttitle=\bfseries,
  title={\stepcounter{thinknum}思考 \thechapter.\thethinknum},
  breakable,
  #1
}

% 动手练习（章末的实操题目）
\newtcolorbox{exercise}[1][]{
  colback=cyan!3,
  colframe=cyan!40!black,
  fonttitle=\bfseries,
  title={\stepcounter{exercisenum}练习 \thechapter.\theexercisenum},
  breakable,
  #1
}
